introduccion
La adaptación de textos es un proceso de transformación lingüística que busca ajustar la dificultad de un contenido escrito para satisfacer las necesidades de un público objetivo específico. En el ámbito del procesamiento de lenguaje natural (NLP), esta tarea se aborda tradicionalmente desde dos vertientes: la simplificación de textos (Text Simplification, TS), que reduce la complejidad léxica y sintáctica, y la complejización, que busca enriquecer el texto para desafiar a lectores con competencias más avanzadas.

Aunque la adaptación es bidireccional, la vasta mayoría de la literatura y los esfuerzos de investigación se han concentrado en la simplificación. Históricamente, el objetivo principal de la TS ha sido mejorar la accesibilidad para lectores con bajos niveles de alfabetización o desafíos cognitivos (Siddharthan, 2014). Sin embargo, en el contexto de la enseñanza de lenguas extranjeras, surge un requisito adicional y más restrictivo: la simplificación controlada por niveles de suficiencia (Readability-Controlled Text Simplification o RCTS).

En este trabajo, el enfoque se centra exclusivamente en la simplificación de textos, entendida como la capacidad de transformar un material de entrada de alta complejidad hacia niveles inferiores de la escala CEFR. Esta necesidad ha sido recientemente validada por iniciativas internacionales como la shared task TSAR 2025 (Alva-Manchego et al., 2025), la cual propone simplificar textos de niveles B2 o superiores hacia niveles específicos como A2 o B1. Bajo esta premisa, nuestra investigación no solo busca reducir la complejidad general, sino garantizar que el texto resultante sea pedagógicamente adecuado para el estadio de aprendizaje del estudiante, preservando estrictamente el significado y la coherencia del mensaje original.

trabajo relacionado

diseño de experimentos
prompts
metodos y metricas de evaluacion
analisis resultados
discucion
conclucion
trabajo futuro